\section{Adam --- general MAP minimiser}
\label{sec:adam}

\paragraph{How it works.}
Adam minimises the objective~\eqref{eq:objective} by gradient descent with a per-coordinate
adaptive step, positivity by clamping, a learning rate halved whenever the loss rises, and a
small-change stop:
\begin{quote}\small
$f \leftarrow f_0$ (the ridge $s_2^2$ start);\quad
repeat: $g_k \leftarrow \nabla L(f)$;\ $f \leftarrow f - \text{lr}\cdot \hat m/(\sqrt{\hat v}+\varepsilon)$;\
$f \leftarrow \max(f,0)$;\ every $K$ steps halve lr if the loss rose, stop if it barely moved.
\end{quote}
Because $\lvert \hat m/\sqrt{\hat v}\rvert \approx 1$, each voxel moves by about \texttt{lr} per
iteration \emph{whatever the gradient} --- so \texttt{lr} is a displacement in the units of $f$
(${\sim}0.02$), not a gradient scaling. $L$ is convex, so with a prior that weighs the minimiser
is unique and Adam reaches it; the reconstruction is then only as good as the analytic prior $R$.

\paragraph{Parameters.}
The only axes that move the \emph{solution} are the prior \texttt{reg} and its share $\lambda$.
We compare two priors: \textbf{L1} (sparsity --- few bright voxels) and \textbf{Tikhonov} (the
L2 of the gradient --- smoothness). Tikhonov is a degree-2 prior, so at $f\approx0.02$ it is
${\sim}50\times$ weaker than a degree-1 prior for the same nominal $\lambda$ and over-smooths once
it bites. Everything else is a convergence knob, fixed for the whole campaign: \texttt{lr}$=0.1$
(large enough to move; the scheduler halves it on overshoot), \texttt{init}$=$ridge with weight
$\lambda_{rr}=s_2^2=36.24$ (never the adjoint start, ${\sim}850\times$ too large at MA-TIRF
scale), \texttt{max\_iter}$=5000$ and \texttt{EPS}$=10^{-8}$ so the plateau stop --- not the
budget --- ends the run.

\paragraph{Standard use.}
The smooth prior at a moderate share, $\texttt{reg}=\text{Tikhonov}$, $\lambda=0.1$. It is a safe
default that never empties the image, but it over-smooths depth --- rarely the best prior on the
sparse structures, though closer to competitive on the continuous membrane.

\paragraph{Optimal use.}
Match the prior to the object and raise it with the noise. On the sparse structures (filaments,
vesicles) the sparsity prior L1 has a clear optimum, while the smooth Tikhonov prior only
over-smooths --- its NMSE stays high and flat:
\IfFileExists{tables/adam_shape.tex}{\begin{center}\begin{minipage}{\linewidth}\centering
\captionof{table}{Adam --- prior shape at the reference structure (noise 0.02): NMSE and depth error versus lambda for the sparse (L1) and smooth (Tikhonov) priors.}
\begin{adjustbox}{max width=\linewidth}
\begin{tabular}{llll}
\toprule
prior & lambda & NMSE & depth err (nm) \\
\midrule
l1 & 0.05 & 0.478 & 48 \\
l1 & 0.1 & 0.507 & 54 \\
tikhonov & 0.05 & 0.405 & 66 \\
tikhonov & 0.2 & 0.459 & 80 \\
tikhonov & 0.3 & 0.483 & 87 \\
tikhonov & 0.5 & 0.527 & 97 \\
\bottomrule
\end{tabular}\end{adjustbox}\end{minipage}\end{center}
}%
  {\emph{[pending the run — \texttt{python -m benchmarks analyze}]}}

\noindent L1 minimises around $\lambda \approx 0.2$--$0.3$; Tikhonov, being smooth and degree-2,
cannot localise a filament network and trails it throughout. The prior winner may itself shift with
the object (the smooth prior is less penalised on a continuous membrane) --- the measured
standard-vs-optimal figures, per structure:
\IfFileExists{tables/algo_adam.tex}{\begin{center}\begin{minipage}{\linewidth}\centering
\captionof{table}{ADAM on MA-TIRF --- reconstruction metrics per ground truth and noise level.}
\begin{adjustbox}{max width=\linewidth}
\begin{tabular}{llllllll}
\toprule
dataset & noise & params & NMSE & PSNR & Depth Error (nm) & Stack Recovery & chi2\_ratio \\
\midrule
cell & 0 & reg=tikhonov, lambda=0.1 & 0.72 & 15.2 & 60 & 0 & 0.00507 \\
cell & 0 & reg=none, lambda=0 & 0.476 & 17.0 & 21 & 0 & 1.17e-06 \\
fibres & 0 & reg=tikhonov, lambda=0.1 & 0.321 & 29.0 & 58 & 0.241 & 0.0298 \\
fibres & 0 & reg=none, lambda=0 & 0.0812 & 34.9 & 14 & 0 & 1.59e-05 \\
vesicles & 0 & reg=tikhonov, lambda=0.1 & 0.481 & 29.3 & 69 & 0.5 & 0.048 \\
vesicles & 0 & reg=none, lambda=0 & 0.0697 & 37.7 & 2 & 0 & 1.1e-05 \\
cell & 0.02 & reg=tikhonov, lambda=0.1 & 0.534 & 16.5 & 26 & 0.0758 & 0.653 \\
cell & 0.02 & reg=l1, lambda=0.2 & 0.277 & 19.4 & 16 & 0 & 0.591 \\
fibres & 0.02 & reg=tikhonov, lambda=0.1 & 0.429 & 27.7 & 72 & 0.0431 & 0.433 \\
fibres & 0.02 & reg=tikhonov, lambda=0.05 & 0.405 & 28.0 & 66 & 0.0862 & 0.415 \\
vesicles & 0.02 & reg=tikhonov, lambda=0.1 & 0.552 & 28.7 & 69 & 0 & 0.423 \\
cell & 0.05 & reg=tikhonov, lambda=0.1 & 0.584 & 16.1 & 30 & 0.0606 & 0.615 \\
cell & 0.05 & reg=l1, lambda=0.2 & 0.487 & 16.9 & 22 & 0 & 0.596 \\
fibres & 0.05 & reg=tikhonov, lambda=0.1 & 0.679 & 25.7 & 115 & 0.0431 & 0.411 \\
vesicles & 0.05 & reg=tikhonov, lambda=0.1 & 0.787 & 27.1 & 146 & 0 & 0.391 \\
esoubies & native & reg=l1, lambda=0.2 & — & — & — & — & 18.9 \\
\bottomrule
\end{tabular}\end{adjustbox}\end{minipage}\end{center}
\begin{center}\begin{minipage}{\linewidth}\centering
\captionof{table}{ADAM on deconvolution --- reconstruction metrics per noise level.}
\begin{adjustbox}{max width=\linewidth}
\begin{tabular}{lllllll}
\toprule
dataset & noise & params & NMSE & PSNR & SSIM & chi2\_ratio \\
\midrule
img\_001 & 0 & reg=tikhonov, lambda=0.1 & 0.195 & 14.1 & 0.302 & 0.021 \\
img\_001 & 0 & reg=none, lambda=0 & 0.0191 & 24.4 & 0.694 & 1.96e-07 \\
img\_001 & 0.02 & reg=tikhonov, lambda=0.1 & 0.0293 & 22.6 & 0.576 & 1.09 \\
img\_001 & 0.05 & reg=tikhonov, lambda=0.1 & 0.0371 & 21.5 & 0.511 & 1.09 \\
\bottomrule
\end{tabular}\end{adjustbox}\end{minipage}\end{center}
}{}

\IfFileExists{figures/matirf_fibres_ADAM_noise002.png}{%
\begin{figure}[H]\centering
  \includegraphics[width=.7\linewidth]{matirf_fibres_ADAM_noise002}
  \caption{Adam's best reconstruction of \texttt{fibres} at noise 0.02 (depth map + $yz/zx$
  profiles).}
\end{figure}}{}

\paragraph{When it works, when it fails.}
Adam has no failure of its own on a well-posed problem: it minimises the energy exactly. On
MA-TIRF the quality is decided entirely by $R$ and $\lambda$, and none of the analytic priors
models a cell's depth structure the way a learned denoiser does --- \emph{its ceiling is the
ceiling of hand-built priors}. The failure modes are at the ends of $\lambda$: $\lambda\to0$ is
non-negative least squares (noisy / depth-collapsed), $\lambda\to1$ drops the data (empty or
flat). Adam is the right \emph{engine} and, on MA-TIRF, the wrong \emph{prior}: the gain is in
the regulariser, not the optimiser.
